\begin{theorem}[window-\texorpdfstring{$6$}{6} boundary hidden sector 的地址层与电荷层严格分裂]\label{thm:conclusion-window6-boundary-address-charge-strict-splitting}
\leanverified{paper\_conclusion\_window6\_boundary\_address\_charge\_strict\_splitting}
记
\[
X_6^{\mathrm{bdry}}=\{\texttt{100001},\texttt{100101},\texttt{101001}\},
\qquad
\widetilde X_6^{\mathrm{bdry}}
:=
\bigsqcup_{w\in X_6^{\mathrm{bdry}}}
\bigl(\Fold_6^{\mathrm{bin}}\bigr)^{-1}(w).
\]
则 window-\(6\) 的 boundary hidden sector 具有如下两层彼此不可约化的规范结构：
\begin{enumerate}
  \item \textbf{地址层：}
  \[
  \widetilde X_6^{\mathrm{bdry}}
  \ \text{是一个正则}\ \ZZ_6\text{-torsor},
  \qquad
  \CC[\widetilde X_6^{\mathrm{bdry}}]\cong \CC[\ZZ_6];
  \]
  \item \textbf{电荷层：}
  \[
  P_{\partial}:=Z_6^{\mathrm{bd}}\cong T_6[2]\cong (\ZZ/2\ZZ)^3.
  \]
\end{enumerate}
因此，boundary hidden sector 既不是单纯的三比特立方体，也不是单纯的六点地址 torsor；前者会遗失真正的 \(3\)-周期地址成分，后者则不足以恢复三条中心 parity 轴的独立电荷坐标。
\end{theorem}
\begin{proof}
由定理 \ref{thm:conclusion-window6-boundary-z6-torsor-local-global-mismatch}，\(\widetilde X_6^{\mathrm{bdry}}\) 上存在自由且传递的 \(\ZZ_6\)-作用，故它是一个正则 \(\ZZ_6\)-torsor，其复置换模即为正则表示
\[
\CC[\widetilde X_6^{\mathrm{bdry}}]\cong \CC[\ZZ_6].
\]
另一方面，由定理 \ref{thm:conclusion-window6-boundary-torus-2torsion}，
\[
Z_6^{\mathrm{bd}}\cong T_6[2]\cong (\ZZ/2\ZZ)^3.
\]
前者含有不可约的 \(3\)-周期地址结构，而后者是纯 \(2\)-挠中心电荷立方体；二者在阶、表示型与可见接口上均不可能互相压缩。证毕。
\end{proof}

\begin{corollary}[boundary 标量 parity 的唯一等变--几何重合]\label{cor:conclusion-window6-boundary-unique-equivariant-geometric-scalar}
\leanverified{paper\_conclusion\_window6\_boundary\_unique\_equivariant\_geometric\_scalar}
在 \(P_{\partial}\cong (\FF_2)^3\) 上赋予 boundary 三轴循环置换诱导的 \(C_3\)-作用。则一切 \(C_3\)-等变线性泛函
\[
\ell:P_{\partial}\longrightarrow \FF_2
\]
构成一维 \(\FF_2\)-空间；除零映射外，唯一的 \(C_3\)-等变标量读出为总 parity
\[
\ell_{\mathrm{diag}}(x_1,x_2,x_3)=x_1+x_2+x_3.
\]
并且该唯一等变标量读出恰与强局域 geometric uplift 唯一保留的对角 \(\ZZ/2\) 方向一致。
\end{corollary}
\begin{proof}
由定理 \ref{thm:conclusion-window6-boundary-c3-diagonal-irreducible-splitting}，
\[
P_{\partial}\cong \mathbf 1\oplus V_2,
\]
其中 \(\mathbf 1=\langle(1,1,1)\rangle\) 为唯一非零平凡子模。故
\[
\Hom_{C_3}(P_{\partial},\FF_2)\cong \Hom_{C_3}(\mathbf 1,\FF_2)
\]
是一维空间，其唯一非零元正是总 parity。另一方面，由定理 \ref{thm:conclusion-window6-boundary-parity-zero-one-three-law}，强局域 geometric uplift 只保留
\[
P_6^{\mathrm{geo}}=\langle(1,1,1)\rangle\cong \ZZ/2\ZZ.
\]
两者遂严格重合。证毕。
\end{proof}

\begin{theorem}[boundary 的 \texorpdfstring{$F_8/F_9$}{F8/F9} 双隐通道正交性]\label{thm:conclusion-window6-boundary-f8-f9-channel-orthogonality}
设 \(b_6^{\mathrm{bin}}\) 为首个被截断的 \(F_8\) 修复位，并令 canonical boundary parity 由三条 boundary 二点纤维上的反对称换位定义。则：
\begin{enumerate}
  \item \(b_6^{\mathrm{bin}}\!\restriction_B\equiv 0\)，其中
  \[
  B:=\bigcup_{w\in X_6^{\mathrm{bdry}}}\bigl(\Fold_6^{\mathrm{bin}}\bigr)^{-1}(w);
  \]
  \item canonical boundary parity 在每条 boundary 二点纤维上必取互补值；
  \item 因而任何保持 boundary 反对称性的二值读出都不可能由 \(F_8\) 修复位通道实现，而只能落在 \(F_9\) 的二层别名通道上。
\end{enumerate}
换言之，\(F_8\) 只负责一阶离散修复，而 \(F_9\) 才承载 canonical boundary sheet parity。
\leanverified{paper\_conclusion\_window6\_boundary\_f8\_f9\_channel\_orthogonality}
\end{theorem}
\begin{proof}
由定理 \ref{thm:conclusion-window6-boundary-parity-not-measurable-from-f8}，
\[
b_6^{\mathrm{bin}}\!\restriction_B\equiv 0,
\]
故 \(F_8\) 通道在每条 boundary 二点纤维上恒定。另一方面，由推论 \ref{cor:window6-f8-repair-vs-f9-boundary-parity-separated}，canonical boundary parity 分属 \(F_9\) 的二层别名通道，并在同纤维两点上以反对称方式取值。故两者不可能由同一隐藏通道实现。证毕。
\end{proof}

\begin{theorem}[boundary 的地址刚性与电荷非唯一性]\label{thm:conclusion-window6-boundary-address-rigidity-charge-nonuniqueness}
\leanverified{paper\_conclusion\_window6\_boundary\_address\_rigidity\_charge\_nonuniqueness}
设 \(A\) 为有限阿贝尔群。若存在保持完整 boundary 对称的等变单射
\[
\Phi:\widetilde X_6^{\mathrm{bdry}}\hookrightarrow A,
\qquad
\Phi(g\cdot x)=\rho(g)+\Phi(x)\qquad (g\in \ZZ_6),
\]
则 \(\rho:\ZZ_6\to A\) 必为单射，从而 \(A\) 必含有阶 \(6\) 元；特别地，最小可行有限阿贝尔载体恰为
\[
A_{\min}\cong \ZZ_6\cong \ZZ_2\times \ZZ_3,
\]
并因此满足任何二进制实现都必须有
\[
2^B\ge 6
\qquad\Longrightarrow\qquad
B\ge 3.
\]
与此相对，对每个 boundary parity 赋值
\[
\chi\in \Hom(Z_6^{\mathrm{bd}},\{\pm1\}),
\]
其全局一维特征延拓集恰有
\[
2^{18}
\]
个元素。故 window-\(6\) 的 boundary hidden sector 同时呈现两种刚性完全不同的层：
\begin{enumerate}
  \item 地址层具有 order-\(6\) 载体刚性；
  \item 电荷层则具有 \(2^{18}\)-重全局符号非唯一性。
\end{enumerate}
\end{theorem}
\begin{proof}
若 \(\rho(g)=0\) 对某个 \(0\neq g\in \ZZ_6\) 成立，则
\[
\Phi(g\cdot x)=\Phi(x)
\qquad (\forall x\in \widetilde X_6^{\mathrm{bdry}}).
\]
由 \(\Phi\) 的单射性得 \(g\cdot x=x\) 对所有 \(x\) 成立，这与 \(\widetilde X_6^{\mathrm{bdry}}\) 上 \(\ZZ_6\)-torsor 作用的自由性矛盾。故 \(\rho\) 必为单射，因而 \(A\) 含有一个阶 \(6\) 元。最小此类有限阿贝尔群即为 \(\ZZ_6\)，于是任何二进制载体都须满足 \(2^B\ge 6\)，从而 \(B\ge 3\)。

另一方面，由推论 \ref{cor:conclusion-window6-boundary-conductor-two-plancherel-uniformity}，restriction 到
\[
\Hom(Z_6^{\mathrm{bd}},\{\pm1\})
\]
的纤维大小恰为 \(2^{18}\)。这说明局部 boundary parity 电荷虽然只有三位，但每个赋值都落在一个大小正好为 \(2^{18}\) 的全局符号簇中。故地址层与电荷层分别表现为刚性与非唯一性的严格并存。证毕。
\end{proof}

\begin{theorem}[有理不可见而 torsion 可见的 boundary 超选分裂]\label{thm:conclusion-window6-boundary-rational-blind-torsion-superselection}
\leanverified{paper\_conclusion\_window6\_boundary\_rational\_blind\_torsion\_superselection}
window-\(6\) 的 boundary parity 同时满足：
\begin{enumerate}
  \item 在有理连续包络中完全不可见：对每个 \(q>0\)，一切 \(\QQ\)-系数特征类与有理同伦证书在 \(Z_6^{\mathrm{bd}}\) 上的拉回都为零；
  \item 在复表示层完全可见：任一携带 \(Z_6^{\mathrm{bd}}\)-作用的 fold-compatible 复表示 \(M\) 都分解为
  \[
  M=\bigoplus_{\chi\in \widehat Z_6^{\mathrm{bd}}} M_{\chi},
  \]
  即严格的八重超选分解。
\end{enumerate}
因此，boundary parity 不是新的有理连续零模，而是一个纯 \(2\)-挠的超选阴影：对有理连续证书彻底透明，而对复表示谱分解完全可见。
\end{theorem}
\begin{proof}
第一条由定理 \ref{thm:conclusion-window6-boundary-parity-directsummand-rational-blindness} 给出。第二条由定理 \ref{thm:conclusion-window6-boundary-torus-2torsion} 给出：
\[
Z_6^{\mathrm{bd}}\cong T_6[2]\cong (\ZZ/2\ZZ)^3,
\]
故 Fourier 幂等元分解产生八个角色扇区。两者并置即得“有理不可见而 torsion 可见”的超选分裂。证毕。
\end{proof}

\begin{conjecture}[boundary torsor--charge 分裂原则]\label{conj:conclusion-boundary-torsor-charge-splitting-principle}
设某一更高窗口 \(m\) 的 boundary hidden sector 同时满足：
\begin{enumerate}
  \item 几何 boundary 提升点集形成一个有限 \(\ZZ_{2n}\)-torsor；
  \item 中心 boundary parity 电荷群满足
  \[
  P_{\partial,m}\cong (\ZZ/2\ZZ)^n;
  \]
  \item boundary 轴带有循环 \(C_n\)-作用，且 \(P_{\partial,m}\) 在 \(\FF_2[C_n]\)-模意义下含唯一平凡直线。
\end{enumerate}
则应有如下统一分裂：
\begin{enumerate}
  \item 任何完全对称保持的有限阿贝尔编码载体都必须含有阶 \(2n\) 元；
  \item 强局域 geometric uplift 只保留对角 \(\ZZ_2\subset (\ZZ/2\ZZ)^n\)；
  \item 一切 \(C_n\)-等变标量 parity 读出都压缩为总 parity；
  \item 剩余的 \(n-1\) 位 parity 只能以 torsion 超选数据的形式存活，而不能升级为新的连续有理质量零模。
\end{enumerate}
window-\(6\) 恰是该原则在 \(n=3\) 时的完整实现。
\end{conjecture}
